%% tables/table2_auc_ihca.tex
\begin{table}[H]
\caption{AUC comparison across clinical encoding strategies --- IHC-A arm.
  All models use the DyAM architecture with CT radiomics (PC, PL, LN),
  pathology IHC-A, genomics, and PD-L1~TPS as base modalities;
  clinical encoding is the varied factor.
  AUC computed on pooled out-of-fold predictions from 10-fold
  cross-validation ($n = 247$).
  95\%~CI by DeLong structural components method.
  $\Delta$AUC relative to the no-clinical baseline.
  $^\dagger$Best-performing NLP variant; n/a, not applicable.
  NLP, natural language processing; PCA, principal component analysis.}
\label{tab:auc_ihca}
\begin{tabular}{lccc}
\toprule
\textbf{Model} & \textbf{AUC} & \textbf{95\%~CI} & \boldmath{$\Delta$}\textbf{AUC} \\
\midrule
No clinical (baseline)        & 0.764 & [0.695--0.833] & (ref) \\
+ Labs (13-dim numeric)       & 0.768 & [0.700--0.837] & $+0.004$ \\
+ NLP raw (384-dim)           & 0.781 & [0.717--0.845] & $+0.017$ \\
\textbf{+ NLP-PCA16 (16-dim)}\textsuperscript{$\dagger$}
                              & \textbf{0.784} & \textbf{[0.719--0.848]} & $\mathbf{+0.020}$ \\
+ Labs + NLP (combined)       & 0.767 & [0.700--0.835] & $+0.003$ \\
+ Labs + NLP-PCA16 (combined) & 0.753 & [0.682--0.823] & $-0.011$ \\
+ BioClinBERT-PCA16           & 0.767 & [0.701--0.833] & $+0.003$ \\
\bottomrule
\end{tabular}
\end{table}
